\section{Wprowadzenie do kolekcjonerskich gier karcianych}
Kolekcjonerska gra karciana (ang. \textit{collectible card games} --- CCG) to rodzaj gry turowej, w której faza przygotowań jest równie ważna, jak decyzje podejmowane podczas rozgrywki. W przeciwieństwie do klasycznych gier o pełnej informacji, takich jak szachy czy go, gracze nie rozpoczynają partii z identycznym zestawem bierek lub kart. Zamiast tego przystępują do rozgrywki z samodzielnie skonstruowaną talią (ang. \textit{deck}). Jednymi z najpopularniejszych tego typu gier są Magic: The Gathering \cite{churchill2019} oraz Pokémon \cite{guimaraes2026}. W grach CCG talia tworzona jest z obszernej i stale aktualizowanej puli dostępnych kart. Jak zauważają Yang i in. \cite{yang2021}, proces budowania talii jest nie tylko istotnym elementem gier typu CCG, ale również problemem optymalizacji o rozległej przestrzeni poszukiwań. Potwierdzają to współczesne eksperymenty nad Legends of Code and Magic, gdzie wykazano, że nawet w uproszczonym środowisku gry typu CCG faza budowania talii (ang. \textit{draft phase}) ma kluczowe znaczenie dla sukcesu agenta. Jak opisują Kowalski i Miernik, w nowszych wersjach gry (v1.5) agenci muszą optymalizować talię w czasie rzeczywistym z losowo generowanych kart, co stanowi istotne wyzwanie dla algorytmów \cite{kowalski2023}. Co więcej, aby sztuczna inteligencja mogła w pełni wykorzystać potencjał tak skonstruowanych talii podczas partii, badacze ci stosują techniki ewolucyjne do automatycznego strojenia funkcji oceniających bieżący stan gry \cite{Kowalski2021}.

Wraz z rozwojem technologii narodziła się cyfrowa wersja tego gatunku (ang. \textit{digital collectible card game} --- DCCG). Przeniesienie gry karcianej do świata wirtualnego pozwoliło na wprowadzenie wielu mechanik, które nie byłyby fizycznie możliwe do zrealizowania. Gry DCCG wprowadzają do rozgrywki wyższy stopień stochastyczności i losowości \cite{hoover2020, xue2016}, co --- podobnie jak w przypadku innych gier z niepełną informacją \cite{Luo2026} --- czyni je trudnym środowiskiem testowym dla algorytmów AI \cite{Vieira2022}. Ponadto umożliwiają niestandardowe efekty, takie jak klonowanie kart, generowanie nowych zasobów czy losowe rozdzielanie obrażeń. Jednymi z najpopularniejszych tego typu gier są Gwint: Wiedźmińska gra karciana \cite{kominiarczuk2017}, The Elder Scrolls: Legends \cite{vieira2024} oraz Hearthstone \cite{sanchez2020}.